\subsection{Dataset Details}  \label{subsec:dataset_details}

\begin{table}[t]
    \footnotesize
    \centering
    \setlength{\tabcolsep}{0.3em}
    \begin{tabular}{l c c c c c}
        \toprule
            \textbf{Dataset} & \textbf{Intents \#} & \textbf{Training} & \textbf{Dev} &\textbf{Test} & \textbf{Total} \\
        \midrule
            SNIPS & 7 & 50,625 & 2,613 & 2,615 & 55,853 \\
            ATIS & 18 & 20,250 & 1,125 & 1,125 & 22,500 \\
            Banking77 & 77 & 36,390 & 2,009 & 2,021 & 40,420 \\
            CLINC150 & 147 & 54,896 & 2,889 & 2,977 & 60,762 \\
        \bottomrule
    \end{tabular}
    \caption{Statistics of the constituents of BlendX.}
    \centering
    \label{table:BlendX_statistics}
\end{table}

To summarize this section, we introduce \textbf{BlendX}, a collection of enhanced multi-intent datasets, shaped by our concatenation strategies and further validated using our three newly developed metrics.
We apply our framework to four single-intent datasets: ATIS, SNIPS, Banking77, and CLINC150---resulting in BlendATIS, BlendSNIPS, BlendBanking77, and BlendCLINC150.
The statistics of BlendX are listed in Table \ref{table:BlendX_statistics}.

In the preprocessing step, similar to  MixX \citeplanguageresource{qin-etal-2020-agif}, BlendX ensures that each intent has equal number of utterances in the datasets. 
BlendX may also necessitate duplicating certain utterances, maintaining a ratio of single-, double-, and triple-intent utterances at 3:5:2.
We use single-intent datasets without rectifying their internal errors, and exclude utterances that do not correspond to specific requests.\footnote{e.g., the utterances in CLINC150 whose intents are \texttt{yes}, \texttt{no}, or \texttt{maybe}.} 
For the case of ATIS, each subset of data (training, dev, and test) contains unique types of intents. 
When adapting this to BlendATIS, we retain these characteristics, ensuring it remains relatively more challenging.

We develop the final version of BlendX using both the Manual and Generative concatenation approaches.
For the Manual approach, we keep a balanced ratio for the data instances corresponding to the settings of AND variants, various conjunctions, inherent ambiguities, and gerund phrases, each at 1:1:1:1.
Additionally, using the Generative approach, we create data instances that amount to half of those produced with the AND variants configuration.
They encompass all potential variations of concatenation, with the expectation that they foster natural and intuitive constructions, potentially featuring omissions and coreferences.

\paragraph{Enhancing data quality through data filtering} 
Through the filtering process, we successfully eliminate errors similar to those highlighted in Table \ref{table:generative}. 
This process actively exploits the three metrics proposed in \S\ref{subsec:metric}. 
Initially, we remove clear failures generated by ChatGPT, such as explicit mentions of an intent label or unnecessary punctuation. 

Subsequently, we check all concatenated sentences using the proposed metrics, paying special attention to and filtering out instances where there is a significant discrepancy before ($ \sum_{i=1}^n |utt_i|_x $) and after concatenation ($|utt|_{x}$). 
For example, a substantial increase in word count may indicate unnecessary paraphrasing by ChatGPT, while a significant decrease might suggest overlooked utterances during concatenation. 
Both $\mathrm{C}(utt, n)$ and $\mathrm{P}(utt, n)$ are also subjected to similar filtering logic. 

Lastly, these filtered instances are reviewed by three human experts, with a focus on excluding instances only if they compromise the intended meaning. 
Instances categorized under \textit{intent removed}, \textit{intent changed}, and \textit{failed to merge} were identified and removed due to their significant deviation from the original intent, as detailed in Table \ref{table:generative}. 
This filtering process significantly enhances the quality of our datasets.